% Note PDFs styled like a mathematics textbook (pandoc -H, xelatex); used by pn/notes.py.
% Each note defines \notenumber („Úloha 11/05“) for the running head (empty if it has no number).
\providecommand{\notenumber}{}

% Latin Modern throughout (text, code; maths through unicode-math's default)
\usepackage{fontspec}
\usepackage{xcolor}
\setmainfont{lmroman10}[Extension=.otf, UprightFont=*-regular, BoldFont=*-bold,
  ItalicFont=*-italic, BoldItalicFont=*-bolditalic, SmallCapsFont=lmromancaps10-regular]
\setsansfont{lmsans10}[Extension=.otf, UprightFont=*-regular, BoldFont=*-bold,
  ItalicFont=*-oblique, BoldItalicFont=*-boldoblique]
\setmonofont{lmmono10}[Extension=.otf, UprightFont=*-regular, ItalicFont=*-italic]

% symbols the Latin Modern text fonts lack come from the maths font (also in code)
\usepackage{newunicodechar}
\newunicodechar{⇄}{\ensuremath{\rightleftarrows}}
\newunicodechar{≤}{\ensuremath{\leq}}
\newunicodechar{≥}{\ensuremath{\geq}}
\newunicodechar{≈}{\ensuremath{\approx}}
\newunicodechar{⇒}{\ensuremath{\Rightarrow}}
\newunicodechar{∧}{\ensuremath{\land}}
\newunicodechar{λ}{\ensuremath{\lambda}}
\newunicodechar{μ}{\ensuremath{\mu}}
\newunicodechar{ρ}{\ensuremath{\rho}}
\newunicodechar{τ}{\ensuremath{\tau}}
\newunicodechar{⁶}{\textsuperscript{6}}
\newunicodechar{⁷}{\textsuperscript{7}}
\newunicodechar{⁻}{\textsuperscript{\textminus}}
\newunicodechar{②}{\textcircled{\scriptsize 2}}

\definecolor{accent}{HTML}{1F3A5F}
\definecolor{boxbg}{HTML}{F4F6F9}
\definecolor{rulegray}{gray}{0.55}

% section headings: numbered, small caps, a thin rule under them
\usepackage{titlesec}
\setcounter{secnumdepth}{1}
\titleformat{\section}{\normalfont\Large\scshape\color{accent}}{\thesection}{0.8em}{}[\vspace{2pt}{\color{rulegray}\titlerule[0.4pt]}]
\titlespacing*{\section}{0pt}{2.2ex plus .5ex}{1.4ex}
\titleformat{\subsection}[runin]{\normalfont\bfseries}{}{0pt}{}[.]
\titlespacing*{\subsection}{0pt}{1.6ex plus .4ex}{0.6em}

% lists: compact, round bullets
\usepackage{enumitem}
\setlist{itemsep=0.35ex, topsep=0.6ex, parsep=0pt}
\setlist[itemize,1]{label={\color{accent}\raisebox{0.15ex}{\scriptsize$\bullet$}}, leftmargin=1.4em}
\setlist[itemize,2]{label={\color{accent}\raisebox{0.15ex}{\scriptsize$\circ$}}}

% the task (and the original wording): a rounded box, like an exercise in a textbook
\usepackage[skins,breakable]{tcolorbox}
\renewenvironment{quote}{%
  \begin{tcolorbox}[enhanced, breakable, colback=boxbg, colframe=accent!35, boxrule=0.6pt,
    arc=6pt, left=10pt, right=10pt, top=6pt, bottom=6pt,
    before skip=1.2ex, after skip=1.4ex]\itshape}{\end{tcolorbox}}

% code: small, on a rounded light background, long lines wrap
\usepackage{fvextra}
\DefineVerbatimEnvironment{Highlighting}{Verbatim}{breaklines,breakanywhere,commandchars=\\\{\},fontsize=\footnotesize}
\BeforeBeginEnvironment{Shaded}{\begin{tcolorbox}[enhanced, breakable, colback=boxbg, colframe=boxbg, boxrule=0pt, arc=6pt, left=8pt, right=8pt, top=4pt, bottom=4pt]}
\AfterEndEnvironment{Shaded}{\end{tcolorbox}}

% figures: captioned as in a book; placed where they are
\usepackage{caption}
\captionsetup{font=small, labelfont={sc}, labelsep=period}
\makeatletter\def\fps@figure{!htbp}\makeatother

% a diagram wider than the text column may use the full page width less 2 cm margins,
% centred on the page; narrower diagrams keep their natural size
\makeatletter
\newlength\figmaxwidth
\setlength\figmaxwidth{\dimexpr\paperwidth-4cm\relax}
\renewcommand*\pandocbounded[1]{%
  \sbox\pandoc@box{#1}%
  \Gscale@div\@tempa{\textheight}{\dimexpr\ht\pandoc@box+\dp\pandoc@box\relax}%
  \Gscale@div\@tempb{\figmaxwidth}{\wd\pandoc@box}%
  \ifdim\@tempb\p@<\@tempa\p@\let\@tempa\@tempb\fi
  \makebox[\linewidth][c]{\ifdim\@tempa\p@<\p@\scalebox{\@tempa}{\usebox\pandoc@box}\else\usebox{\pandoc@box}\fi}%
}
\makeatother

% running head in light grey without a rule (a section rule often follows right below)
\usepackage{fancyhdr}
\pagestyle{fancy}
\fancyhf{}
\fancyhead[L]{\small\scshape\color{rulegray} Petriho sítě}
\fancyhead[R]{\small\scshape\color{rulegray}\notenumber}
\fancyfoot[C]{\small\thepage}
\renewcommand{\headrulewidth}{0pt}
\fancypagestyle{plain}{\fancyhf{}\fancyfoot[C]{\small\thepage}\renewcommand{\headrulewidth}{0pt}}

\usepackage{microtype}
\setlength{\parindent}{0pt}
\setlength{\parskip}{0.7ex plus .2ex}
\AtBeginDocument{\hypersetup{colorlinks=true, linkcolor=accent, urlcolor=accent}}
